\chapter{Introduction}
\label{cap:introduction}

\todobox{Contesto di tesi e motivazione applicativa}{%
Inquadrare l'attivit\`a svolta nell'ambito della tesi di laurea magistrale:
\begin{itemize}
  \item Collaborazione con il gruppo di ricerca di riferimento (Politecnico di Torino).
  \item Esigenze dell'agricoltura di precisione: diserbo meccanico selettivo, monitoraggio delle colture e riduzione dell'impatto ambientale senza dipendenza da segnali GNSS RTK sotto chioma o in serra.
  \item Obiettivo di realizzare una soluzione autonoma, modulare e fruibile su calcolatori embedded a basso costo (famiglia NVIDIA Jetson).
\end{itemize}
}

\section{Context and motivation}
\label{sec:context}

Global population growth and rising food demand exert increasing pressure on agricultural productivity, while the sector faces persistent labor shortages and rising operating costs~\cite{bai2023vision,silva2023rowdetection}. Addressing these demands requires higher automation in field operations such as weeding, spraying, and crop monitoring. Row-based crops are among the most studied applications of agricultural robots and account for about 75\% of the planted cropland in the United States~\cite{navone2025gpsfree}. Autonomous ground vehicles designed for these settings must navigate along inter-row corridors reliably without damaging crops or compacting productive soil beds.

Outdoor agricultural navigation historically relied on Global Navigation Satellite Systems (GNSS) augmented by Real-Time Kinematic (RTK) corrections. However, satellite positioning presents well-documented limitations in field environments: thick crop canopies cause multi-path reflections and severe signal attenuation, while greenhouse installations obstruct satellite reception entirely~\cite{bai2023vision,navone2025gpsfree}. Furthermore, RTK base station infrastructure adds cost. Alternative exteroceptive sensors such as 3D LiDAR provide direct spatial measurements, but their high unit cost hinders economic scalability across commercial robotic fleets, and their point clouds lack the color and textural information required to distinguish crops from weeds~\cite{bai2023vision,navone2025gpsfree}. Monocular RGB cameras offer a cost-effective alternative that captures detailed chromatic and texture cues directly from the scene.

Early vision-based guidance pipelines extracted navigation lines through handcrafted image processing techniques. These approaches compute color vegetation indices to isolate plant matter from soil, apply intensity thresholding such as Otsu's method, and detect corridor edges using line-fitting algorithms or Hough transforms~\cite{bai2023vision,mendez2023vineyards}. Although computationally lightweight, these calibrated pipelines degrade when ambient field conditions shift. Variable sunlight angles, cast shadows from foliage, soil moisture gradients, specular reflections on wet leaves, and weed infestations frequently disrupt fixed thresholding rules and generate false boundary estimates.

Deep learning architectures, particularly convolutional neural networks for semantic segmentation, overcome these hand-tuning limitations. By learning hierarchical visual representations directly from image datasets, segmentation networks achieve higher tolerance to illumination shifts, ground irregularities, and weed clutter than handcrafted color models.

However, deep neural networks introduce a practical operational bottleneck: supervised training requires extensive pixel-level ground truth. Manually labeling agricultural scenes is labor-intensive and must cover multiple crop types, growth stages, and ambient conditions. In addition, running neural networks on compact onboard computers requires balancing segmentation accuracy against processing speed and power consumption, as steering commands must be updated fast enough to guide the vehicle in real time.

This thesis addresses these challenges through a modular pipeline that combines automatic dataset annotation with SAM~3, real-time onboard semantic segmentation, and deterministic geometric reference extraction for local robot guidance.

\section{Objectives of the thesis}
\label{sec:objectives}

The main objectives of this work are:
\begin{itemize}
\item Evaluate whether segmentation networks trained on SAM~3 labels, used without manual correction, approach a network trained on manual masks, using manually annotated lettuce images as ground truth.
\item Select a segmentation architecture balancing accuracy and inference speed for robotic deployment, through offline evaluation of FCN, DeepLabV3, and LR-ASPP on manually annotated lettuce images.
\item Evaluate the effect of output label configurations (four-class, three-class, and binary path) on path segmentation accuracy.
\item Implement an edge-based geometric extractor that converts segmentation masks into local image-space navigation references, and compare it with three baseline algorithms.
\item Integrate the selected perception and geometric components on a robot platform for field validation.
\end{itemize}

\section{Organization of the work}
\label{sec:organization}

The thesis is organized as follows:

\textbf{Chapter~\ref{cap:background}} introduces the agricultural navigation problem, the role of vision-based local guidance, and the segmentation-to-control pipeline structure.

\textbf{Chapter~\ref{cap:stato_arte}} reviews vision-based methods for agricultural row navigation, covering classical and learning-based segmentation approaches, geometric extraction techniques, and experimental validation in related work.

\textbf{Chapter~\ref{cap:architectures}} details the mathematical formulations and designs of the deep learning architectures evaluated in this work, focusing on FCN, DeepLabV3, MobileNetV3 with the LR-ASPP head, and SAM~3.

\textbf{Chapter~\ref{cap:segmentation}} describes dataset preparation with manual annotation and automatic SAM~3 labeling, empirical comparisons among the three architectures on lettuce, and evaluations of supervision strategies across label configurations.

\textbf{Chapter~\ref{cap:geometric_reference}} presents the geometric extractor that converts corridor masks into navigation references through horizontal-strip interval selection and independent edge fitting, together with three baseline algorithms used for comparison.

\textbf{Chapter~\ref{cap:platform}} describes the robot hardware, camera configuration, software integration using ROS, and the deployment of the trained segmentation model.

\textbf{Chapter~\ref{cap:experiments}} reports offline and field validation results for segmentation accuracy, geometric reference extraction, and closed-loop navigation performance.

\textbf{Chapter~\ref{cap:conclusions}} summarizes the main findings, discusses limitations, and identifies directions for future work.
